\documentclass[10pt,a4paper]{amsart}
\usepackage[margin=19mm]{geometry}
\usepackage{amsmath,amssymb,mathtools}
\usepackage[scr=rsfs,cal=euler]{mathalfa}
\usepackage{xfrac}
\usepackage[unicode,hidelinks,bookmarks=false]{hyperref}
\newcommand{\ie}{i.e.\ }
\newcommand{\cf}{cf.\ }
\newcommand{\resp}{respectively\ }
\newcommand{\Subsection}{\subsection}
\providecommand{\nobreakdash}{\nobreak\hbox{-}}
\setlength{\emergencystretch}{2em}
\multlinegap0pt
\usepackage{bm}
\usepackage{bbm}
\usepackage{dsfont}
\usepackage{xspace}
\usepackage{cancel}
\usepackage{mathtools}
\usepackage{amsthm}
\makeatletter
\@ifundefined{lem}{\newtheorem{lem}{Lemma}}{}
\makeatother
\DeclareMathOperator{\Aut}{Aut}
\DeclareMathOperator{\id}{id}
\DeclareMathOperator{\sw}{sw}
\newcommand{\Z}{\mathbb{Z}}
\newcommand{\aesf}{\psi}
\newcommand{\aeso}{\phi}
\newcommand{\aes}{\omega}
\newcommand{\cX}{\mathcal{X}}
\newcommand{\cY}{\mathcal{Y}}
\newcommand{\sr}{ r } %{(0,0,2\alpha_1)}
\title[Sets of equiangular lines in dimension $18$ constructed from]{Sets of equiangular lines in dimension $18$ constructed from $A\_9 \oplus A\_9 \oplus A\_1$}
\author{Yen-chi Roger Lin, Akihiro Munemasa, Tetsuji Taniguchi,, Kiyoto Yoshino}
\date{Source: arXiv:2503.06377v3 (CC BY 4.0)}
\begin{document}
\maketitle

\noindent\emph{Source and licence.} Sets of equiangular lines in dimension $18$ constructed from $A_9 \oplus A_9 \oplus A_1$. Version arXiv:2503.06377v3 (CC BY 4.0), distributed under the Creative Commons Attribution 4.0 International licence, as recorded in the arXiv licence metadata for this exact version and shown on the abstract page https://arxiv.org/abs/2503.06377v3. Full rights notice: licenses/arxiv-2503.06377-CC-BY-4.0.txt.\par\smallskip

\noindent\textit{Editorial scope.} Counted unit: the Lem whose source label is \texttt{lem:genL} in the article (source0.tex lines 863--865, source label \texttt{lem:genL}), delivered together with the article's own complete original proof of it (source0.tex lines 866--934, one \texttt{proof} environment of 69 lines). No mathematics is written, repaired or re-derived: statement and proof are the article's own text line for line. Required context retained uncounted: 1d (lines 387--389); lem:chi (lines 564--569); lem:X5 (lines 667--670); lem:L5 (lines 690--696); lem:X1Y10 (lines 749--760); lem:con (lines 792--804); lem:con:Epm (lines 795--798); lem:con:0 (lines 800--802). Every definition, display and label of those lines is retained unchanged. Omitted from this package and not redistributed: 1--386, 390--563, 570--666, 671--689, 697--748, 761--791, 799--799, 803--862, 935--1368. Nothing inside a retained excerpt is omitted or altered: every retained body is delivered exactly as the article's lines are written, so no omission receipt of any kind accompanies this package. Citations: the retained excerpts carry no citation command at all, so no bibliography entry of the article is transcribed and no reference is redistributed. Reference adaptation: none was needed and none was made; every reference inside the delivered unit resolves inside it, so no unresolved reference or citation is produced. Printed slips kept literal: no printed word, symbol, hypothesis or number of the counted statement or of its proof was altered. Layout and class: the article's own class is \texttt{amsart} with packages including amsmath, amsthm, amsfonts, color, comment, fontenc, times, mathtools, hyperref, enumitem; the wrapper is the lane's pinned \texttt{amsart} editorial wrapper at 10pt on A4, which restates the theorem-like environments the retained text needs and adds the article's own macro definitions that the retained text needs (\texttt{\textbackslash Aut}, \texttt{\textbackslash Z}, \texttt{\textbackslash aes}, \texttt{\textbackslash aesf}, \texttt{\textbackslash aeso}, \texttt{\textbackslash cX}, \texttt{\textbackslash cY}, \texttt{\textbackslash id}, \texttt{\textbackslash sr}, \texttt{\textbackslash sw}), with extra wrapper packages (bm, bbm, dsfont, xspace, cancel, mathtools). The delivered numbering is the unit's own. Assets: no retained span contains a figure environment, a graphics-inclusion command, a PDF-page inclusion, a table or a TikZ picture; no image file is copied into this package, the package declares no asset dependency and no image file is redistributed with it. Licence: version arXiv:2503.06377v3 of this article is distributed under the Creative Commons Attribution 4.0 International licence, as recorded in the arXiv licence metadata for this exact version and shown on the abstract page https://arxiv.org/abs/2503.06377v3. A full rights notice is delivered at licenses/arxiv-2503.06377-CC-BY-4.0.txt. Characters: every retained excerpt is pure ASCII, so the delivered engine, XeLaTeX with its default Latin Modern text font, reports no missing character.
	\begin{align}
		\Lambda / L \simeq \Z/10\Z. \label{1d}
	\end{align}

\begin{lem}	\label{lem:chi}
	We have
	$\cX/\sim_{\sw}
		= \left\{ [\aesf \cup \aeso] : [\aesf] \in \cX_5/\sim_{\sw}, \aeso \in \cX_1 \right\}.$
	In particular, $\left|\cX/\sim_{\sw} \right| = \left|\cX_5/\sim_{\sw} \right| \left| \cX_1 \right|.$
\end{lem}

\begin{lem}
	The number $|\cX_5/\sim_{\sw}|$ of maximum affine equiangular sets in $X_5$ with respect to $\sr$ up to switching equivalence is $151200$.	\label{lem:X5}
	Also, the automorphism group $\Aut(\Lambda)_{\sr}$ transitively acts on $\cX_5/\sim_{\sw}$.
\end{lem}

\begin{lem}	\label{lem:L5}
	We have
	\begin{align*}
		L_5 & = \langle (e_i-e_j,0,0) : i,j \in [4] \rangle \oplus \langle (e_i-e_j ,0,0): i,j \in [10]\setminus[4] \rangle \oplus \langle \sr \rangle \\
		    & \simeq A_3 \oplus A_5 \oplus A_1.
	\end{align*}
\end{lem}

\begin{lem}	\label{lem:X1Y10}
	The maximum affine equiangular sets $\aeso \in \cX_1$ in $X_1$ with respect to $\sr$ correspond one-to-one to the elements $(Y_i)_{i=1}^{10} \in \cY_{5}$ in such a way that
	\begin{align}	\label{lem:X1Y10:0}
		( \{ \{j,k\} \in E(K_{10}) : v_1(i,j,k) \in \aeso \} )_{i=1}^{10}  % \quad \text{ for each } \quad \aes \in \cX_1,
	\end{align}
	in $\cY_5$, and
	\begin{align}	\label{lem:X1Y10:1}
		\left\{ v_1(i,j,k) : \{j,k\} \in Y_i \ \text{ for }\ i \in [10] \right\}%  \quad \text{ for each } \quad (Y_i)_{i=1}^{10} \in \cY_{5}.
	\end{align}
	in $\cX_1$.
	In particular, a maximum affine equiangular set in $X_1$ with respect to $\sr$ is of cardinality $45$.
\end{lem}

\begin{lem}	\label{lem:con}
	Let $(Y_i)_{i=1}^{10} \in \cY_5$.
	Let
	\begin{align}
		E_+ := \bigcup_{i=1}^4 Y_i \quad \text{and} \quad E_- := \bigcup_{i=5}^{10} Y_i.
		\label{lem:con:Epm}
	\end{align}
	Let $G$ be a graph with vertex set $[10]$ and edges $\{j,j'\} \in \binom{[10]}{2}$ satisfying for some $k \in [10]$,
	\begin{align}	\label{lem:con:0}
		\{\{j,k\}, \{j',k\}\} \subset E_+ \ \text{ or }\ \{\{j,k\}, \{j',k\}\} \subset E_-.
	\end{align}
	Then, $G$ is connected.
\end{lem}

\begin{lem}	\label{lem:genL}
	Every maximum affine equiangular set in $X$ with respect to $\sr$ and the switching root $\sr$ generates $\Lambda$.
\end{lem}

\begin{proof}
	Let $\aes' \in \cX$.
	Let $M'$ be the lattice generated by $\aes'$ together with $r$.
	By Lemma~\ref{lem:chi}, we may assume without loss of generality that
	$\aes' = \aesf' \cup \aeso'$ for some
	% there exist 
	$\aesf' \in \cX_5$ and $\aeso' \in \cX_1$.
	By Lemma~\ref{lem:X5}, there exists $\sigma \in S_{10}\times\{\id \}\times\{ \id \}\subset
		\Aut(\Lambda)$ such that
	\begin{align}	\label{lem:genL:1}
		\aesf := \sigma(\aesf') = \psi_0.
	\end{align}
	Let $M:=\sigma(M')$. Then,
	Lemma~\ref{lem:L5} together with~\eqref{lem:genL:1} implies
	\begin{align} \label{lem:genL:3}
		(A_3 \oplus A_5 ) \oplus O_{10} \oplus A_1  \subset M.
	\end{align}
	Here, denote by $O_n$ the lattice $\{ 0 \} \subset \Z^n$ for a positive integer $n$.
	Let $\aeso := \sigma(\aeso')$ and $\aes:=\aesf\cup\aeso$.
	Then, $\aes$ together with $r$ generates $M$.
	Moreover, since $\sigma$ leaves $X_1$ invariant, we have $\aeso \in \cX_1$ .
	%	by Lemma~\ref{lem:aut}.
	By Lemma~\ref{lem:X1Y10}, there exists $(Y_i)_{i=1}^{10} \in \cY_{5}$ such that
	\begin{align} \label{lem:genL:4}
		\aeso = \left\{ v_1(i,j,k) : \{j,k\} \in Y_i \ \text{ for }\  i \in [10] \right\} \subset X_1.
	\end{align}
	Let $E_+$ and $E_-$ be as in~\eqref{lem:con:Epm}.
	%	$E_+ := \bigcup_{i=1}^{4} Y_i $ and $E_- := \bigcup_{i=5}^{10} Y_i$ as in~\eqref{lem:con:-1}.

	First, we prove
	\begin{align} \label{lem:genL:2}
		O_{10} \oplus A_9 \oplus A_1 \subset M.
	\end{align}
	Indeed, we define a graph $G$ with vertex set $[10]$ and edges $\{j, j'\} \in \binom{[10]}{2}$ if
	\begin{align}	\label{lem:genL:6}
		(0,e_{j'}-e_j,0) \in M.
	\end{align}
	It suffices to prove that $G$ is connected.
	Note that, if $\{j,k\}\in Y_i$ and $\{j',k\}\in Y_{i'}$, then
	%	 for $i,i', j \neq k \neq j' \in [10]$,
	\begin{align}	\label{lem:genL:5}
		v_1(i,j,k)-v_1(i',j',k) = (e_{i'}-e_i, e_{j'}-e_j , 0) \in M.
	\end{align}
	Therefore, we have from~\eqref{lem:genL:3},
	$(0,e_{j'}-e_j,0) \in M$, or equivalently, $\{j,j'\} \in E(G)$,
	provided that \eqref{lem:con:0} holds.
	By applying Lemma~\ref{lem:con} to $G$, we conclude that $G$ is connected, that is, \eqref{lem:genL:2} holds.

	Next, we prove
	\begin{align}	\label{lem:genL:7}
		A_9 \oplus A_9 \oplus A_1  \subset M.
	\end{align}
	There exist $i \in [4]$ and $i' \in [10] \setminus [4]$ such that $Y_{i'}$ and $Y_i$ are non-empty.
	Fix $\{ j',k' \} \in Y_{i'}$ and $\{j,k\} \in Y_{i}$.
	Then,
	\begin{align*}
		v_1(i,j,k)-v_1(i',j',k') = (e_{i'} - e_i , e_{j'} + e_{k'} - e_j - e_k , 0) \in M.
	\end{align*}
	By~\eqref{lem:genL:2}, we have
	$
		(e_{i'} - e_i , 0 , 0) \in M.
	$
	By~\eqref{lem:genL:3} and~\eqref{lem:genL:2}, we obtain~\eqref{lem:genL:7}.

	Finally, we prove $M=\Lambda$.
	The elements $\aes \cap X_5$ are of order $2$ in $\Lambda/L$, and the elements $\aes \cap X_1$ are of order $5$ in $\Lambda/L$.
	Since $L \subset M$ by~\eqref{lem:genL:7} and $|\Lambda/L|=10$ by~\eqref{1d},
	we conclude $M= \Lambda$.
\end{proof}

\end{document}
